\documentclass[11pt]{article}
\usepackage[margin=1in]{geometry}
\usepackage{amsmath,amssymb,amsthm}
\usepackage{graphicx}
\usepackage{hyperref}

\newtheorem*{theorem}{Counterexample}
\newtheorem*{lemma}{Lemma}

\title{Duplicate synthesis atoms refute the operator-valued Feichtinger conjectures}
\author{Run fa\_banach\_001, lane 6}
\date{2026-06-19}

\begin{document}
\maketitle

\section*{Source conjectures}

The source paper is K. Mahesh Krishna and P. Sam Johnson,
\emph{Operator-Valued p-Approximate Schauder Frames}, arXiv:2201.03955.
Conjectures 7.1 and 7.2 appear in Section 7, on page 24 of the source PDF.

\begin{center}
\includegraphics[width=\linewidth]{figures/open_problem_crop.png}
\end{center}

Conjecture 7.1 asks whether every operator-valued $p$-approximate Schauder
frame $(\{A_n\}_n,\{\Psi_n\}_n)$ with
\[
0<\inf_n\|\Psi_n\|\leq \sup_n\|\Psi_n\|<\infty,
\qquad
0<\inf_n\|A_n\|\leq \sup_n\|A_n\|<\infty
\]
can be partitioned into finitely many operator-valued $p$-approximate Riesz
sequences.  Conjecture 7.2 states the same assertion for operator-valued
$p$-approximate Bessel sequences.

\section*{Relevant source definitions}

For Banach spaces $X,Y$, an operator-valued $p$-approximate Schauder frame
($p$-ASF) in $\mathcal B(X,Y)$ is a pair
$(\{A_n\}_n,\{\Psi_n\}_n)$ with $A_n\in\mathcal B(X,Y)$ and
$\Psi_n\in\mathcal B(Y,X)$ such that the analysis map
$\theta_Ax=\sum_n L_nA_nx$, the synthesis map
$\theta_\Psi z=\sum_n\Psi_n\Gamma_nz$, and the frame operator
$S_{A,\Psi}=\theta_\Psi\theta_A=\sum_n\Psi_nA_n$ are well defined and
bounded, with $S_{A,\Psi}$ invertible.  If invertibility of $S_{A,\Psi}$ is
not required, the source calls the pair an operator-valued $p$-approximate
Bessel sequence ($p$-ABS).

The source defines an operator-valued $p$-approximate Riesz basis ($p$-ARB) by
the condition
\[
P_{A,\Psi}:=\theta_A S_{A,\Psi}^{-1}\theta_\Psi
   = I_{\ell^p(\mathbb N)\otimes Y}.
\]
A subcollection is an operator-valued $p$-approximate Riesz sequence if it is a
$p$-ARB on
\[
Z=\overline{\operatorname{span}}\bigcup_n \Psi_n(Y).
\]

\section*{Proof Intuition}

The conjectures control only the norms of the individual analysis and
synthesis atoms.  They do not forbid infinitely many different analysis atoms
from sharing exactly the same synthesis atom.  In the example below, all bad
atoms synthesize into the first coordinate vector of $\ell_1$, while their
analysis functionals read all coordinates.  Adding the ordinary coordinate
atoms for the tail makes the total frame operator invertible, so the pair is a
genuine operator-valued $1$-ASF.  But every finite partition puts two of the
duplicate synthesis atoms into the same part, and a Riesz subcollection cannot
have non-injective synthesis.

\begin{lemma}
If a subcollection is an operator-valued $p$-approximate Riesz sequence, then
its synthesis operator is injective.
\end{lemma}

\begin{proof}
For a $p$-ARB, the source definition gives
$I=\theta_A S_{A,\Psi}^{-1}\theta_\Psi$ on the corresponding coefficient
space.  If $\theta_\Psi z=0$, then
\[
z=\theta_A S_{A,\Psi}^{-1}\theta_\Psi z=0.
\]
Thus $\theta_\Psi$ is injective.  The same argument applies to a Riesz
sequence after restricting to its closed span $Z$.
\end{proof}

\begin{theorem}
Conjectures 7.1 and 7.2 of arXiv:2201.03955 are false as stated.  There is an
operator-valued $1$-ASF in $\mathcal B(\ell_1,\mathbb K)$ whose analysis and
synthesis atoms all have norm $1$, but which cannot be partitioned into
finitely many operator-valued $1$-approximate Riesz sequences.  Consequently it
is also a counterexample to the $1$-ABS conjecture.
\end{theorem}

\begin{proof}
Let $X=\ell_1$ over the scalar field $\mathbb K$, let $Y=\mathbb K$, and write
$(e_k)_{k\geq1}$ for the standard unit vector basis of $\ell_1$.  Let
$e_k^*$ be the $k$th coordinate functional, and let
$J_k:\mathbb K\to\ell_1$ be $J_k(\alpha)=\alpha e_k$.

Index the family by the countable set
\[
I=\{b_k:k\geq 1\}\cup\{g_k:k\geq 2\},
\]
identified with $\mathbb N$ by any enumeration.  Define
\[
A_{b_k}=e_k^*,\qquad \Psi_{b_k}=J_1\quad(k\geq1),
\]
and
\[
A_{g_k}=e_k^*,\qquad \Psi_{g_k}=J_k\quad(k\geq2).
\]
All these operators have norm $1$, so the lower and upper norm bounds in
Conjectures 7.1 and 7.2 are satisfied.

Since $Y=\mathbb K$, the coefficient space
$\ell^1(I)\otimes Y$ is canonically $\ell^1(I)$.  For $x=(x_k)\in\ell_1$,
the analysis sequence is
\[
\theta_Ax=\big((x_k)_{k\geq1},(x_k)_{k\geq2}\big)\in
\ell^1(\{b_k\})\oplus_1\ell^1(\{g_k\}),
\]
hence
\[
\|\theta_Ax\|_1=\sum_{k\geq1}|x_k|+\sum_{k\geq2}|x_k|\leq 2\|x\|_1.
\]
For coefficients $a=(a_{b_k},a_{g_k})\in\ell^1(I)$,
\[
\theta_\Psi a=\left(\sum_{k\geq1}a_{b_k}\right)e_1+
 \sum_{k\geq2}a_{g_k}e_k,
\]
so
\[
\|\theta_\Psi a\|_1
\leq \sum_{k\geq1}|a_{b_k}|+\sum_{k\geq2}|a_{g_k}|=\|a\|_1.
\]
Thus the analysis and synthesis maps are bounded.

The frame operator is
\[
Sx=\theta_\Psi\theta_Ax
 =\left(\sum_{k\geq1}x_k\right)e_1+\sum_{k\geq2}x_ke_k.
\]
It is bounded on $\ell_1$.  It is also invertible: for
$y=(y_k)\in\ell_1$,
\[
S^{-1}y=\left(y_1-\sum_{k\geq2}y_k,\; y_2,\; y_3,\ldots\right),
\]
and
\[
\|S^{-1}y\|_1
\leq |y_1|+2\sum_{k\geq2}|y_k|
\leq 2\|y\|_1.
\]
Therefore $(\{A_i\}_{i\in I},\{\Psi_i\}_{i\in I})$ is an operator-valued
$1$-ASF in $\mathcal B(\ell_1,\mathbb K)$.

Now take any finite partition $I=I_1\cup\cdots\cup I_m$.  The bad set
$B=\{b_k:k\geq1\}$ is infinite, so some part $I_j$ contains two distinct bad
indices, say $b_r$ and $b_s$.  In the coefficient space of that subcollection,
the vector $\delta_{b_r}-\delta_{b_s}$ is nonzero, but its synthesis is
\[
\theta_\Psi^{I_j}(\delta_{b_r}-\delta_{b_s})
  =J_1(1)-J_1(1)=0.
\]
Hence the synthesis operator of the $I_j$ subcollection is not injective.  By
the lemma, that subcollection cannot be an operator-valued $1$-approximate
Riesz sequence.  Since this happens for every finite partition, the frame
cannot be partitioned into finitely many operator-valued $1$-approximate Riesz
sequences.

This disproves Conjecture 7.1.  Since every $1$-ASF is, after dropping the
invertibility requirement, a $1$-ABS in the sense of the source definition, the
same example also disproves Conjecture 7.2.
\end{proof}

\section*{Verification notes}

The construction uses only the source paper's definitions.  The key check is
that the duplicated synthesis maps do not prevent boundedness of the synthesis
operator when $p=1$: summation of the bad coefficients is a bounded functional
on $\ell^1$.  The added good coordinate atoms make the total frame operator
invertible by an explicit bounded inverse.

No numerical computation is needed.  The crop was generated by
\texttt{code/make\_open\_problem\_crop.py} from the local source PDF.

\section*{Novelty check and limitations}

A bounded novelty check was performed on 2026-06-19.  The run indexes were
searched for arXiv:2201.03955 and the core phrases ``operator-valued
p-approximate Schauder frame'', ``operator-valued p-ASF'', ``Feichtinger
conjecture'', and ``p-approximate Riesz sequence''.  External web/arXiv
queries for the exact title and for combinations of ``operator-valued
p-approximate Schauder frame'', ``Feichtinger conjecture'', and
``counterexample'' found the source paper and adjacent frame-conjecture papers,
but no direct counterexample to Conjectures 7.1--7.2.

The counterexample is for $p=1$.  It refutes the conjectures as written for
arbitrary $p\in[1,\infty)$, but it does not address possible repaired
statements for fixed $p>1$, for Hilbert/G-frame subcases, or under hypotheses
that prevent repeated synthesis atoms.

\section*{Human-review recommendation}

Review as a likely valid counterexample to Conjectures 7.1 and 7.2 as stated.
The main points to check are the interpretation of a Riesz subcollection on
the closed span $Z$ and the injectivity consequence of the source definition
$P_{A,\Psi}=I$.

\begin{thebibliography}{9}
\bibitem{KrishnaJohnson2022}
K. Mahesh Krishna and P. Sam Johnson,
\emph{Operator-Valued p-Approximate Schauder Frames},
arXiv:2201.03955, 2022.
\end{thebibliography}

\end{document}
